\scanpage{59}
\uline{Graphes cubiques}

Un graphe $\Delta$ est dit \inserted{\uline{graphe}} \uline{cubique} si \deleted{\unsure{les arêtes}} ses arêtes sont déterminées par leurs extrémités, (qui sont distinctes, et s'il existe \deleted{deux} parties $I$, \inserted{$I'$,} de l'ens des sommets satisfaisant les conditions suivantes

a) \deleted{\unsure{$I$ de}} $\operatorname{card} I=\operatorname{card} I'=6$, $I$ et $I'$ disjoints

b) $\forall i\in I,\ \exists!\ i'\in I'$ tel que $i'$ non lié à $I$, et $\forall i'\in I',\ \exists!\ i\in I$ tel que $i$ non lié à $I$ [donc $I$ et $I'$ sont en corr. biunivoque par la relation ``$i$ n. lié à $i'$'' \deleted{\unsure{k}}, soit $\unsure{i|\iota}:\deleted{\unsure{I\to I'}}\ I\cup I'\longrightarrow I\cup I'$] \deleted{\unsure{tel}}

c) $\forall$ partie $J$ \inserted{$=\{i,j\}$} à deux éléments de $I$, $\exists!\ \zeta_J=\zeta_{ij}\in\Delta$ \inserted{$-J\cup i(J)$} qui soit lié aux \unsure{4|a|\deleted{deux}} éléments \deleted{$i,j$ et} de $J\cup i.(J)$, et non lié aux autres éléments de $I\cup I'$.

\deleted{[\unsure{donc c'est donc que} $\lambda_J\notin I\cup I'$}

[Il est clair que $J\longmapsto\lambda_J$ est une appl. \uline{injective} de $\mathfrak P_2(I)$ dans $\Delta-(I\cup I')$

d) L'application précédente
\[
\mathfrak P_2(I)\longrightarrow\Delta-(I\cup I')
\]
est \deleted{bijective} surjective (donc bijective), i.e. $\forall\lambda\in\Delta-I\cup I'$, \deleted{$\exists J\in\mathfrak P_2(\Delta)$ telle que $\lambda$ soit lié à $J\cup i(J)$,} l'ens. des éléments de $I\cup I'$ auxquels $\lambda$ est lié est de cardinal $4$ et stable par $i$.

e) Si $J,J'\in\mathfrak P_2(I)$, $J\ne J'$, alors $\zeta_J$ et $\zeta_{J'}$ sont liés ssi $J\cap J'=\varnothing$.

\centerline{---}

Graphe cubique $\Delta(I)$ défini par ens $I$ de card. $6$.

$\Delta$ cubique ssi $\Delta\simeq\Delta(I)$ (pour un $I$ de card. $6$, ou pour tous)
% archive sheet number: 58
